\documentclass[10pt,a4paper]{amsart}
\usepackage[margin=19mm]{geometry}
\usepackage{amsmath,amssymb,mathtools}
\usepackage[scr=rsfs,cal=euler]{mathalfa}
\usepackage{xfrac}
\usepackage[unicode,hidelinks,bookmarks=false]{hyperref}
\newcommand{\ie}{i.e.\ }
\newcommand{\cf}{cf.\ }
\newcommand{\resp}{respectively\ }
\newcommand{\Subsection}{\subsection}
\providecommand{\nobreakdash}{\nobreak\hbox{-}}
\setlength{\emergencystretch}{2em}
\multlinegap0pt
\usepackage{amssymb}
\usepackage{enumerate}
\newtheorem{corollary}{Corollary}
\title{A new collection of $n-$tuple operator inequalities}
\author{Zameddin I. Ismailov, Pembe Ipek Al, Hamid Reza Moradi,, Mohammad Sababheh}
\date{Source: arXiv:2602.14463v1 (CC BY 4.0)}
\begin{document}
\maketitle

\noindent\emph{Source and licence.} A new collection of $n-$tuple operator inequalities. Version arXiv:2602.14463v1 (CC BY 4.0), distributed by arXiv under the Creative Commons Attribution 4.0 International licence, http://creativecommons.org/licenses/by/4.0/, as recorded in the arXiv licence metadata for this exact version and shown on the abstract page https://arxiv.org/abs/2602.14463v1. Full licence notice: \texttt{licenses/arxiv-2602.14463-CC-BY-4.0.txt}.\par\smallskip

\noindent\textit{Editorial scope.} Counted unit: the article's corollary Cor\_Block\_two (source0.tex lines 418--432). The article's complete original proof of it is delivered with it (source0.tex lines 433--483, one \texttt{proof} environment of 51 lines). No mathematics is written, repaired or re-derived: the statement and the proof are the article's own lines, in the article's own notation, and no retained line is substituted or re-flowed.  Retained uncounted as the source-local context the proof needs: lines 83--88; lines 378--381.  Nothing was omitted from any retained span, so the package carries no omission record.  No retained line contains a citation, so the wrapper carries no bibliography and no bibliography entry.  No retained line contains a figure environment, an image-inclusion command or drawing code, so no image is delivered and the package declares no asset dependency.  Editorial formatting only, in addition: an AMS \texttt{amsart} preamble that restates the article's own declarations and macros used by the retained lines, namely the environments \texttt{corollary} on separate counters, together with the standard packages those lines need.  The retained environments print in this document as Corollary 1 in order of appearance, whereas the article prints the same items with the numbering of its own full document.  The complete native source of the article as distributed in the arXiv e-print for this exact version is bundled unchanged as \texttt{sources/194709/source0.tex}.
\begin{equation}\label{Eq_Ident_Bl}
\omega \left( \left[ \begin{matrix}
   O & {{T}_{1}}  \\
   T_{2}^{*} & O  \\
\end{matrix} \right] \right)=\sup_{\theta\in\mathbb{R}}\|T_1+e^{i\theta}T_2\|,
\end{equation}

\begin{corollary}\label{4}
Let ${{T}_{1}},{{T}_{2}}\in \mathbb B\left( \mathbb H \right)$. Then 
	\[{{\left\| {{T}_{1}}+{{T}_{2}} \right\|}^{2}}\le \left\| \frac{3}{2}\left( {{\left| {{T}_{1}} \right|}^{2}}+{{\left| {{T}_{2}} \right|}^{2}} \right)+\Re\left( T_{1}^{*}{{T}_{2}} \right) \right\|.\]
\end{corollary}

\begin{corollary}\label{Cor_Block_two}
Let $T,{{T}_{1}},{{T}_{2}}\in \mathbb B\left( \mathbb H \right)$. Then 
	\[{{\omega }^{2}}\left( \left[ \begin{matrix}
   O & {{T}_{1}}  \\
   T_{2}^{*} & O  \\
\end{matrix} \right] \right)\le \frac{1}{2}\omega \left( \begin{matrix}
   O & \frac{3}{2}\left( {{\left| {{T}_{1}} \right|}^{2}}+{{\left| {{T}_{2}} \right|}^{2}} \right)  \\
   T_{2}^{*}{{T}_{1}} & O  \\
\end{matrix} \right).\]
In particular,
\[{{\omega }^{2}}\left( {{T}} \right)\le \frac{1}{2}\omega \left( \begin{matrix}
   O & \frac{3}{2}\left( {{\left| {{T}} \right|}^{2}}+{{\left| T^{*} \right|}^{2}} \right)  \\
   T^{2} & O  \\
\end{matrix} \right).\]
\end{corollary}

\begin{proof}
It follows from Corollary \ref{4} that
	\[\begin{aligned}
   {{\left\| {{T}_{1}}+{{T}_{2}} \right\|}^{2}}&\le \left\| \frac{3}{2}\left( {{\left| {{T}_{1}} \right|}^{2}}+{{\left| {{T}_{2}} \right|}^{2}} \right)+\Re\left( T_{1}^{*}{{T}_{2}} \right) \right\| \\ 
 & =\left\| \Re\left( \frac{3}{2}\left( {{\left| {{T}_{1}} \right|}^{2}}+{{\left| {{T}_{2}} \right|}^{2}} \right)+T_{1}^{*}{{T}_{2}} \right) \right\| \\ 
 & \le \omega \left( \frac{3}{2}\left( {{\left| {{T}_{1}} \right|}^{2}}+{{\left| {{T}_{2}} \right|}^{2}} \right)+T_{1}^{*}{{T}_{2}} \right).  
\end{aligned}\]
Indeed, we have shown that
\begin{equation}\label{5}
{{\left\| {{T}_{1}}+{{T}_{2}} \right\|}^{2}}\le \omega \left( \frac{3}{2}\left( {{\left| {{T}_{1}} \right|}^{2}}+{{\left| {{T}_{2}} \right|}^{2}} \right)+T_{1}^{*}{{T}_{2}} \right).
\end{equation}
If we replace ${{T}_{2}}$ by ${{e}^{i\theta }}{{T}_{2}}$, in \eqref{5}, we get
	\[\begin{aligned}
   \frac{1}{4}{{\left\| {{T}_{1}}+{{e}^{i\theta }}{{T}_{2}} \right\|}^{2}}&\le \frac{1}{4}\omega \left( \frac{3}{2}\left( {{\left| {{T}_{1}} \right|}^{2}}+{{\left| {{T}_{2}} \right|}^{2}} \right)+{{e}^{i\theta }}T_{1}^{*}{{T}_{2}} \right) \\ 
 & \le \frac{1}{4}\left\| \frac{3}{2}\left( {{\left| {{T}_{1}} \right|}^{2}}+{{\left| {{T}_{2}} \right|}^{2}} \right)+{{e}^{i\theta }}T_{1}^{*}{{T}_{2}} \right\| \\ 
 & \le \frac{1}{2}\omega \left( \begin{matrix}
   O & \frac{3}{2}\left( {{\left| {{T}_{1}} \right|}^{2}}+{{\left| {{T}_{2}} \right|}^{2}} \right)  \\
   T_{2}^{*}{{T}_{1}} & O  \\
\end{matrix} \right).  
\end{aligned}\]
That is,
	\[\frac{1}{4}{{\left\| {{T}_{1}}+{{e}^{i\theta }}{{T}_{2}} \right\|}^{2}}\le \frac{1}{2}\omega \left( \begin{matrix}
   O & \frac{3}{2}\left( {{\left| {{T}_{1}} \right|}^{2}}+{{\left| {{T}_{2}} \right|}^{2}} \right)  \\
   T_{2}^{*}{{T}_{1}} & O  \\
\end{matrix} \right).\]
Now, by taking the supremum over $\theta \in \mathbb{R}$, \eqref{Eq_Ident_Bl} implies
	\[{{\omega }^{2}}\left( \left[ \begin{matrix}
   O & {{T}_{1}}  \\
   T_{2}^{*} & O  \\
\end{matrix} \right] \right)\le \frac{1}{2}\omega \left( \begin{matrix}
   O & \frac{3}{2}\left( {{\left| {{T}_{1}} \right|}^{2}}+{{\left| {{T}_{2}} \right|}^{2}} \right)  \\
   T_{2}^{*}{{T}_{1}} & O  \\
\end{matrix} \right).\]
In particular, if we replace $T_{2}^{*}$ by ${{T}_{1}}$, we get
	\[\begin{aligned}
   {{\omega }^{2}}\left( {{T}_{1}} \right)&={{\omega }^{2}}\left( \left[ \begin{matrix}
   O & {{T}_{1}}  \\
   {{T}_{1}} & O  \\
\end{matrix} \right] \right) \\ 
 & \le \frac{1}{2}\omega \left( \begin{matrix}
   O & \frac{3}{2}\left( {{\left| {{T}_{1}} \right|}^{2}}+{{\left| T_{1}^{*} \right|}^{2}} \right)  \\
   T_{1}^{2} & O  \\
\end{matrix} \right).  
\end{aligned}\]
Namely,
	\[{{\omega }^{2}}\left( {{T}_{1}} \right)\le \frac{1}{2}\omega \left( \begin{matrix}
   O & \frac{3}{2}\left( {{\left| {{T}_{1}} \right|}^{2}}+{{\left| T_{1}^{*} \right|}^{2}} \right)  \\
   T_{1}^{2} & O  \\
\end{matrix} \right).\]
We deduce the desired result by replacing $T_1$ by $T$.
\end{proof}	

\end{document}
